\begin{problem}[第34届物理决赛][引力红移与等效原理]{86}
    爱因斯坦等效原理可表述为：在有引力作用的情况下的物理规律和没有引力但有适当加速度的参考系中的物理规律是相同的。作为一个例子，考察下面两种情况：
    \begin{subproblem}
        \begin{subsubproblem}
            \begin{subsubsubproblem}
                \begin{subsubsubsubproblem}
                    当一束光从引力势比较低的地方传播到引力势比较高的地方时，其波长变长，这个现象称为引力红移。如果在某质量分布均匀的球形星体表面附近的 $A$ 处竖直向上发射波长为 $\lambda_{0}$ 的光，在 $A$ 处竖直上方高度为 $L$ 的 $B$ 处放置一固定接收器，求 $B$ 处接收器接收到的光的波长 $\lambda'$。已知该星体质量为 $M$，半径为 $R$ ($R\gg L$)；引力场满足弱场条件，可应用牛顿引力理论；真空中的光速为 $c$，引力常量为 $G$。
                \end{subsubsubsubproblem}
                \begin{subsubsubsubproblem}
                    如图，假设在没有引力的情况下，有一个长度为 $L$ 的箱子，在箱子上、下面的 $B$、$A$ 两处分别放置激光接收器和激光发射器。在 $t=0$ 时刻，箱子从初速度为零开始，沿 $AB$ 方向做加速度大小为 $a$ 的匀加速运动 ($aL\ll c^{2}$)，同时从 $A$ 处发出一束波长为 $\lambda_{0}$ 的激光。根据狭义相对论，求 $B$ 处接收器接收到的激光的波长 $\lambda''$。
                \end{subsubsubsubproblem}
                \begin{subsubsubsubproblem}
                    根据等效原理，试比较（i）和（ii）的结果，要使物理规律在（i）和（ii）中的情况下相同，则（ii）中的 $a$ 应为多大？
                \end{subsubsubsubproblem}
            \end{subsubsubproblem}
        \end{subsubproblem}
        \begin{subsubproblem}
            引力红移现象的第一个实验验证是在地球表面附近利用穆斯堡尔探测器完成的，穆斯堡尔探测器能以极高的精度分辨伽马光子的能量。按第（1）（i）问，在地球表面附近，$A$ 处放置一个静止的伽马辐射源，辐射的伽马光子的频率为 $\nu_{0}$；$B$ 处放置一个穆斯堡尔探测器，假设该探测器在相对于自身静止的参考系中仅能探测到频率为 $\nu_{0}$ 的伽马光子。为了探测到从 $A$ 处发射的伽马光子，该穆斯堡尔探测器需要某一竖直向下的运动速度。1960-1964年期间，庞德、雷布卡和斯奈德利用美国哈佛大学杰弗逊物理实验室的高塔多次做了这个实验，实验中 $L=22.6\,\mathrm{m}$。试问：$A$ 处发射的伽马光子被探测到时，该穆斯堡尔探测器的运动速度为多大？已知地球表面重力加速度 $g=9.81\,\mathrm{m\cdot s^{-2}}$，真空中的光速 $c=3.00\times10^{8}\,\mathrm{m\cdot s^{-1}}$。
        \end{subsubproblem}
    \end{subproblem}
\end{problem}